% !TeX program = xelatex
% ============================================================================
%  第 12 课　费马小定理的证明　—— 教师版讲义
%  与 02-课件/第12课-费马小定理/第12课-费马小定理.tex 同步
% ============================================================================

\jhlesson{12}{费马小定理的证明}{时长 45 分钟\quad·\quad 教具：无}

\begin{mubiao}
  \begin{itemize}
    \item 能说清为什么两边可以划掉相同的部分。
    \item 能从具体数字过渡到用字母写的一般形式。
    \item 能写出 \(a^{\,p-1}\equiv1\pmod p\)，并知道它叫费马小定理。
  \end{itemize}
\end{mubiao}

\noindent{\small 本课节奏：0—6 回顾｜6—26 排列相乘｜26—38 划掉、再换成字母｜38—45 收尾。}

\section*{一、回顾（0—6 分钟）}

把上节课的猜想原文写在黑板正中：\(p\) 素数、\(a\) 不是 \(p\) 的倍数时 \(a^{\,p-1}\equiv1\pmod p\)。

强调一句：上节课试过的 \(10\) 组都对，但「试出来的」不算数。今天证明它。

\section*{二、排列相乘（6—26 分钟）}

取 \(p=7\)、\(a=3\)，把 \(1,2,3,4,5,6\) 每个都乘 \(3\)，两行写成一张表：

\begin{center}
\begin{tabular}{@{}lcccccc@{}}
  \toprule
  原来的数 & \(1\) & \(2\) & \(3\) & \(4\) & \(5\) & \(6\) \\
  \midrule
  乘 \(3\) 之后 & \(3\) & \(6\) & \(9\) & \(12\) & \(15\) & \(18\) \\
  取余数 & \(3\) & \(6\) & \(2\) & \(5\) & \(1\) & \(4\) \\
  \bottomrule
\end{tabular}
\end{center}

\begin{caozuo}
  先写前两行，余数一行让学生报。报完把两行念一遍：
  上排 \(1,2,3,4,5,6\)，下排 \(3,6,2,5,1,4\)。
  让学生自己说出那句话：\textbf{下排是上排的一次重排}，一个不多、一个不少。
\end{caozuo}

这一步其实第 3 课就见过——「乘 \(3\)」把 \(1\) 到 \(6\) 六个位置重新排了一遍。

然后把两行各自乘起来：

\[
  1\times2\times3\times4\times5\times6=6!, \qquad
  3\times6\times2\times5\times1\times4=6! .
\]

因为下排只是上排的重排，两个乘积一模一样。

\begin{zhu}
  第一次出现的 \(6!\) 要停下来讲一句：\(1\times2\times3\times4\times5\times6\) 写起来太长，
  数学里缩写成 \(6!\)，读作「\(6\) 的阶乘」。学生版讲义里也注了同一句话，不用他们自己记。
\end{zhu}

换个角度再看下排：它正是 \(3\times1,\;3\times2,\;\dots,\;3\times6\) 的余数，所以

\[
  3\times6\times9\times12\times15\times18
  =3^{6}\times(1\times2\times3\times4\times5\times6)
  =3^{6}\times6! ,
\]

两边取模 \(7\)：

\[
  3^{6}\times6!\equiv6!\pmod 7 .
\]

\section*{三、划掉，再换成字母（26—38 分钟）}

抛问题：两边都有的 \(6!=1\times2\times3\times4\times5\times6\)，能不能划掉？

答案：\textbf{能}。模 \(7\) 里，只要不是 \(7\) 的倍数就一定有逆元（第 3 课学的）。
\(6!\) 的因数里没有一个 \(7\)，所以它可逆，两边能同时划掉。于是

\[
  3^{6}\equiv1\pmod 7 .
\]

\begin{caozuo}
  手指回答：我们把 \(6!\) 划掉，靠的是它的什么？（靠它有逆元）
  \(3^6\) 到底等于多少？除以 \(7\) 余几？（\(729=104\times7+1\)，余 \(1\)）
\end{caozuo}

\begin{zhu}
  「为什么能划掉」是本课最关键的一步，也是学生最容易含糊过去的一步。
  一定要让他们说出「因为它有逆元」，而不是「因为两边一样」。
\end{zhu}

现在把 \(7\) 换成 \(p\)、把 \(3\) 换成 \(a\)，左边具体、右边一般，一栏一栏对着写：

\begin{center}
\begin{tabular}{@{}ll@{}}
  \toprule
  \textbf{具体}（\(p=7\)，\(a=3\)） & \textbf{一般} \\
  \midrule
  \(3,\ 6,\ 9,\ \dots,\ 18\) & \(a,\ 2a,\ 3a,\ \dots,\ (p-1)a\) \\
  余数是 \(1\text{—}6\) 的重排 & 余数是 \(1\text{—}(p-1)\) 的重排 \\
  \(3^{6}\times6!\equiv6!\) & \(a^{\,p-1}\times(p-1)!\equiv(p-1)!\) \\
  \(3^{6}\equiv1 \pmod 7\) & \(a^{\,p-1}\equiv1 \pmod p\) \\
  \bottomrule
\end{tabular}
\end{center}

一句话：从数字换成字母，就是「\textbf{抽象}」这两个字的实际意思。

把整个证明收成两块砖：

\begin{itemize}
  \item \textbf{砖一}：\(a,2a,\dots,(p-1)a\) 的余数，正好是 \(1,2,\dots,p-1\) 的一次重排；
  \item \textbf{砖二}：模 \(p\) 里只要不是 \(p\) 的倍数就有逆元，可以两边同时划掉。
\end{itemize}

\begin{dingli}[费马小定理]
  设 \(p\) 是素数，\(a\) 是不被 \(p\) 整除的整数，那么
  \[
    a^{\,p-1}\equiv1 \pmod p .
  \]
\end{dingli}

\section*{四、收尾（38—45 分钟）}

把定理原文读一遍。然后指回黑板最右边——第 1 课写下的那行 \(999999\div7=?\)。

\begin{zhu}
  那行字整门课都没擦，今天还是不要回答。只说一句：「下节课，我们去回答它。」
\end{zhu}

\begin{zhu}
  下课前一分钟把两件事说清楚：（1）下节课要算身份证上的一串数字；
  （2）请他们把号码抄在讲义上带来（没有身份证的同学，抄课后问来的一串 \(18\) 位号码即可）。
\end{zhu}

\begin{xiaojie}
  \begin{itemize}
    \item 证明第一步：乘 \(a\) 取余，只是把 \(1,\dots,p-1\) 重排了一遍。
    \item 两边相乘，再划掉都有的 \((p-1)!\)（因为它有逆元）。
    \item 得到费马小定理：\(p\) 素数、\(a\) 不是 \(p\) 的倍数时，\(a^{\,p-1}\equiv1\pmod p\)。
  \end{itemize}
\end{xiaojie}

\noindent\textbf{下节课}：它有什么用——先去回答黑板右边那个问题。\\
\textbf{下节课要带}：学生自己的身份证号码（或一串 \(18\) 位号码）。